\section{模型的检验与灵敏度分析}

模型检验由完整约束审计和参数灵敏度实验共同完成。前者检查任务唯一执行、实时任务即时开工、到达与截止、单向时延、GPU 容量、IT 功率、设施功率、第 $2406$ 小时禁占用、能源守恒、SOC 递推、终端 SOC、充放电互斥和购售电边界。实际运行的约束满足状态为真，且全时域可行标记为真，说明输出轨迹在所采用数值容差和附件边界下保持内部一致。该检验属于代理验证和物理一致性检验，不使用不可观测的最优标签。

灵敏度实验以问题二系统运行成本 $-419658147.7075976$ 元为基准，分别改变购电价格水平和新能源出力水平。实验直接复算扰动后的目标值及相对变化，不从曲线外推未测试的参数区间。

\subsection{购电价格灵敏度}

购电价格系数变化下的运行成本见表~\ref{tab:price_sensitivity}，对应曲线见图~\ref{fig:sensitivity_price}。

\begin{table}[htbp]
\centering
\caption{购电价格变化对系统运行成本的影响}
\label{tab:price_sensitivity}
\begin{tabular}{rrrr}
\toprule
参数变化 & 参数系数 & 运行成本（元） & 相对变化（\%）\\
\midrule
$-20\%$ & $0.8$ & $-419660613.261396$ & $-0.0005875148169636924$\\
$-10\%$ & $0.9$ & $-419659380.48449683$ & $-0.00029375740848894776$\\
$10\%$ & $1.1$ & $-419656914.9306984$ & $0.00029375740848894776$\\
$20\%$ & $1.2$ & $-419655682.15379924$ & $0.0005875148169636924$\\
\bottomrule
\end{tabular}
\end{table}

表~\ref{tab:price_sensitivity} 显示，购电价格下降 $20\%$ 时，运行成本相对基准变化 $-0.0005875148169636924\%$；价格上升 $20\%$ 时，相对变化为 $0.0005875148169636924\%$。当扰动缩小至 $10\%$ 时，目标变化的绝对值相应为 $0.00029375740848894776\%$。因此，在本次购售电结构和固定调度方案下，目标函数对整体购电价格水平表现出较强稳健性。

\begin{figure}[htbp]
\centering
\includegraphics[width=0.82\textwidth]{sensitivity_electricity_price.png}
\caption{系统运行成本对购电价格水平的灵敏度}
\label{fig:sensitivity_price}
\end{figure}

图~\ref{fig:sensitivity_price} 呈现了价格系数变化与运行成本之间近似对称的响应方向，价格提高使目标值上升，价格降低使目标值下降。曲线在已测试区间内的纵向变化相对基准成本较小，与表~\ref{tab:price_sensitivity} 中较低的相对变化幅度一致。该稳健性只适用于整体价格同比例缩放，不能直接推广至峰谷价差、售电价格或任务调度随价格联动变化的情形。

\subsection{新能源水平灵敏度}

新能源出力水平变化下的运行成本见表~\ref{tab:renewable_sensitivity}，对应曲线见图~\ref{fig:sensitivity_renewable}。

\begin{table}[htbp]
\centering
\caption{新能源水平变化对系统运行成本的影响}
\label{tab:renewable_sensitivity}
\begin{tabular}{rrrr}
\toprule
参数变化 & 参数系数 & 运行成本（元） & 相对变化（\%）\\
\midrule
$-20\%$ & $0.8$ & $-302819505.6070327$ & $27.84138536063257$\\
$-10\%$ & $0.9$ & $-378911868.9949424$ & $9.709397740812053$\\
$10\%$ & $1.1$ & $-445329416.4918755$ & $-6.1171858391189184$\\
$20\%$ & $1.2$ & $-457431461.6153355$ & $-9.000972366216745$\\
\bottomrule
\end{tabular}
\end{table}

新能源水平下降 $20\%$ 时，运行成本相对基准变化 $27.84138536063257\%$；下降 $10\%$ 时，相对变化仍达到 $9.709397740812053\%$。新能源水平上升 $10\%$ 和 $20\%$ 时，运行成本相对变化分别为 $-6.1171858391189184\%$ 和 $-9.000972366216745\%$。由此可得，模型目标对新能源水平较为敏感，且新能源下降造成的成本不利变化在已测试区间内更为显著。

\begin{figure}[htbp]
\centering
\includegraphics[width=0.82\textwidth]{sensitivity_renewable_level.png}
\caption{系统运行成本对新能源出力水平的灵敏度}
\label{fig:sensitivity_renewable}
\end{figure}

图~\ref{fig:sensitivity_renewable} 表明，随着新能源水平提高，运行成本持续下降，说明新能源直接消纳与外送收益对目标具有重要影响。曲线响应并不呈现与价格实验相同的小幅对称变化，反映了新能源不足时购电缺口扩大与新能源富余时外送受限之间的非线性差异。因而，在模型推广和运行决策中，新能源预测误差应被视为比整体购电价格同比例误差更重要的不确定性来源。

综合两组实验可知，系统运行成本对购电价格整体缩放表现稳健，而对新能源水平变化较为敏感。该结论建立在任务方案和能源分流规则保持一致的条件下；由于实际实现没有针对每个扰动重新执行完整任务邻域与储能动作搜索，灵敏度结果衡量的是固定启发式结构下的参数响应，而不是重新优化后的价值函数灵敏度。
